%% Conclusion.
%% Included by sigconf.tex. Do not add a preamble here.

\section{Conclusion}
Irregularly sampled data is the product of a temporal process and an observation process, and the second can shift
independently of the first. We formalised that shift, gave a controlled protocol that re-observes held-out
histories while preserving their targets, and found degradation ranging from a few percent to more than an order of
magnitude across seven task settings, in a way that observation density alone does not explain. The method we
propose exploits the one asymmetry the setting offers: during training, and only during training, the same history
can be observed more richly than deployment allows. Privileged observation self-distillation turns that richer
observation into a task-level target for a student restricted to the deployment view, and it improves absolute
accuracy where the sparse view leaves the target ambiguous. Observation-process augmentation supplies the student's
views and is what reduces sensitivity to the shift itself; hard observation-process mining improves worst-case
error, but only when its candidate pool is harsher than the mixture augmentation already draws, which is why we
report it as a conditional extension rather than a component. None of the three changes what deployment requires:
one sparse observation and the student. The natural next questions are how the gain varies with the size of the
teacher--student observation gap, and whether a shift between two real acquisition environments behaves like the
synthetic re-observations studied here.
